\section{Related work}
\label{sec:related}

We screened roughly two hundred benchmarks and methods. Table~\ref{tab:related-full} (Appendix~\ref{app:related}) sets the closest against the five properties that this benchmark combines: a fresh instance on every exposure, recurrence after variable gaps, a task set that grows through the stream, binary feedback that is charged, and a protocol open to prompting, memory and weight updates. We found no work that has all five.

\paragraph{ARC.}
\citet{chollet2019measure} notes that ARC's format is close-ended, a weakness that motivates this work, and suggests that a better one would let the test-taker query an example generator and submit answers for feedback repeatedly, with the feedback consumed as the score. To our knowledge, this has not been implemented, although Re-ARC \citep{hodel2024rearc} supplies the generators and verifiers it needs. The strongest ARC-AGI-1 systems instead adapt their weights at test time and discard the adaptation, per task \citep{akyurek2024ttt,franzen2025poe} or over the evaluation set \citep{pourcel2025soar}, never meeting a fresh instance of a task they have adapted to. ArcMemo \citep{ho2025arcmemo} writes concepts to memory, but over a single fixed pass. Over a stream of generated ARC-style tasks, \citet{zhang2026faulty} find that a memory the model keeps rewriting first helps and then degrades. ARC-AGI-3 \citep{arcprize2026arcagi3} moves to learning over time within a single game, and a harness that keeps per-game notes across runs \citep{bedi2026learningmachines} has saturated its public games, untested on others.

\paragraph{Continual learning.}
The classic protocol \citep{lopezpaz2017gradient,vandeven2019three} is a fixed sequence of tasks with free labels, scored on fixed test sets after each task boundary. Where tasks come back at all, they return as a planned second pass \citep{wolczyk2021continual,powers2022cora,veniat2021efficient} or as unannounced revisits \citep{caccia2020online,hemati2023class}. Fresh instances are standard in reinforcement learning environments \citep{cobbe2020procgen,kuttler2020nethack} and in generated reasoning suites with verifiers \citep{stojanovski2025reasoninggym,chen2025enigmata,zeng2026rlve}, but these suites train first and test afterwards, and supervision is free.

\paragraph{Agent benchmarks.}
StreamBench \citep{wu2024streambench} gives binary feedback in one shuffled pass, MemoryBench \citep{ai2025memorybench} feedback from simulated users, LifelongAgentBench \citep{zheng2025lifelongagentbench} brings skills back but not instances, and SWE-Bench-CL \citep{joshi2025swebenchcl} regroups SWE-Bench by repository. CL-Bench \citep{asawa2026clbench} scores a stateful system against a stateless one, and AgentCL \citep{shu2026agentcl} and PATH-Bench \citep{yang2026pathbench} rerun a stream or reinsert a probe task. All draw from a fixed pool of instances that appear once or recur unchanged. None grows the task set, charges for supervision or serves many users, and most evaluate prompting and memory only. Concurrent work \citep{anonymous2026continualautoresearch,zheng2026servelearnbench,zhao2026cle} neither charges for feedback nor trains weights.

\paragraph{Memory against weights.}
The comparison between memory and weights exists only in pieces. Memory-based agents \citep{shinn2023reflexion,zhao2024expel,suzgun2025dynamiccheatsheet,zhou2025memento,nekoei2025jefhinter,zhang2026ace} are motivated by the cost or forgetting of fine-tuning but seldom compare against it, and the closest to our program library keeps a skill as a program once it runs correctly \citep{wang2025asi}. On-policy reinforcement learning is reported to forget less than supervised fine-tuning \citep{shenfeld2025rlsrazor,chen2026retaining}, but to forget severely over a sequence of tasks \citep{luo2026rlforgets}, as does harness learning as tasks drift \citep{chen2026selfevolving}. SEAL \citep{zweiger2025seal} runs self-directed weight updates on ARC, test-time curricula \citep{hubotter2025testtimecurricula} train on tasks chosen for the target, and self-improving harnesses \citep{karten2026primeagent} rewrite their memory and skills, but each is evaluated on its own terms. Concurrently, Self-Direction \citep{kansal2026selfdirection} distills an agent's reflections into a LoRA adapter, on task sets recurring at most three times. Memory and weights have yet to be compared on tasks that keep returning.
